

\section{Convergence of Convolution Powers}
\label{sec:3.2}
% BEGIN SUBJECT INDEX: sec:3.2
\index{convolution powers!convergence}
\index{kernel!limiting measure}
% END SUBJECT INDEX


Start with definition of topological semigroup.  
Given countable discrete example.

Now consider the case when $(X_k)$ is an i.i.d.\ sequence of random variables.  
Thus $\mu_j = \mu$ (say) for all $j$.

     Suppose
\[
     \mu^n \to \nu \hbox{ as } n \to \i.
\]
(This is what we mean by stochastic stability.)
Then $\mu * \mu^n = \mu^{n+1}$, letting $n \to \i$
\[
\begin{aligned}
       &\Rightarrow \quad \mu * \nu = \nu                  \\
       &\qquad \qquad \hbox{using the continuity of convolution} \\
       &\Rightarrow \quad \mu^k * \nu = \nu \hbox{ for all } k         \\
       &\Rightarrow \quad \nu * \nu = \nu.
\end{aligned}
\]
Thus $\nu$ is an idempotent measure.

Thus idempotents will play a key role in our study of the convergence of convolution powers of a probability measure.

One of the most important concepts is that of the support of a
measure.

\begin{definition}   \label{defn:3.2}
Let $\mu \in \Prob(S)$.  
Then $S(\mu) = \{ y \in S \mid \mu(y) > 0 \}$.  
$S(\mu)$ is called the support of $\mu$.
\end{definition}
 
\begin{proposition} 	\label{prp:3.1}
Let $\mu,\nu \in \Prob(S)$.  
Then
\[
     S(\mu * \nu) = S(\mu) \cdot S(\nu).
\]
\end{proposition}

\begin{proof}  
Let $y \in S(\mu * \nu)$
\[
     \Rightarrow \quad \mu * \nu(y) > 0
\]
but
\[
     \mu * \nu(y) = \sum_{s_1s_2 = y} \mu(s_1) \nu(s_2) > 0
\]

\[
\begin{aligned}
       &\Rightarrow \quad \F s_1,s_2 \hbox{ with } \mu(s_1) \nu(s_2) > 0 \\
       &\Rightarrow \quad s_1 \in S(\mu), \ s_2 \in S(\nu)             \\
       &\Rightarrow \quad y = s_1 \cdot s_2 \in S(\mu) \cdot S(\nu).
\end{aligned}
\]
\qed
\end{proof}

\runinhead{Homework.} Prove the converse.

Notice the above proposition.  
Show that if $\mu^2 = \mu$, $S(\mu) = S(\mu * \mu) = S(\mu) \cdot S(\mu)$.  
Thus the support of an idempotent measure is a semigroup.

\runinhead{Question.} What kind of semigroup??

Notice the following.  
Let $f: s \to \BDR$ be any real-valued function defined on $S$ and suppose $\nu = \nu * \mu$.  
Then for each $z \in S$,
\[
     \nu(z) = \sum_{st=z} \nu(s) \mu(t).
\]
Thus
\[
\begin{aligned}
       f(z) \nu(z) &= f(z) \sum_{st=z} \nu(s) \mu(t)           \\
       \sum_{t\in S} f(z) \nu(z) &= \sum_{z\in S}
               \bigg[ \sum_{st=z} f(z) \nu(s) \mu(t) \bigg]    \\
       &= \sum_{z\in S}
               \bigg[ \sum_{st=z} f(st) \nu(s) \mu(t) \bigg]
\end{aligned}
\]
let
\[
\begin{gathered}
     G(z) = \{ (s,t) \mid st = z \} \ub S \times S     \\
     \bud_{z\in S} G(z) = S \times S.
\end{gathered}
\]
Notice $G(z_1) \cap G(z_2) = \emptyset$, $z_1 \not= z_2$.  But
$(s_1,s_2) \in S \times S$ then if $s_1s_2 = s$, $(s_1,s_2) \in G(s)$.
Clearly, $\bud G(z) \ub S \times S$.  Thus,
\[
\begin{aligned}
       \sum_{z\in S} f(z) \nu(z) &= \sum_{z\in S}
               \bigg[ \sum_{(s,t)\in G(z)} f(st)
                         \nu(s) \mu(t) \bigg]                  \\
       &= \sum_{S\in S} \sum_{t\in S} f(st) \nu(s) \mu(t).
\end{aligned}
\]
In fact, this is how convolution is defined in the general case via
the Riesz representation theorem.

\begin{lemma} 		\label{lem:3.1} 
Let $\mu,\nu \in \Prob(S)$ such that $\nu = \nu * \mu = \mu * \nu$.  
Then for any $B \ub S$, $x \in S(\nu)$, $y \in S(\mu)$,
\[
\begin{aligned}
       \nu(Bx^{-1}y^{-1}) &= \nu(Bx^{-1})          \\
       \nu(y^{-1}x^{-1}B) &= \nu(x^{-1}B).
\end{aligned}
\]
\end{lemma}

\begin{proof} Let $x \in S(\nu)$, $B \ub S$
\[
\begin{aligned}
       &\nu(Bx^{-1}) = a                           \\
       &A = \{ s \in S \mid \nu(B s^{-1}) \le a \}.
\end{aligned}
\]
Now $x \in A \Rightarrow \nu(A) > 0$.

     Define $g(s) = \max (\nu(Bs^{-1}) - a, 0)$.  Then $\sum\limits_t
\bigg[ \sum\limits_s g(st) \mu(s) - g(t) \bigg] \nu(t) = 0$.  Now
\[
\begin{aligned}
       \nu(Bt^{-1}) - 1 &= \sum_s \big[
               \nu(Bt^{-1}s^{-1}) - a \big] \mu(s)             \\
       &= \sum_j \nu [B(st)^{-1} - a] \mu(s)                   \\
       &\le \sum_s g(st) \mu(s)
\end{aligned}
\]
but
\[
     \sum_s g(st) \mu(s) \ge 0.
\]
Thus
\[
\begin{gathered}
       \max \big( \nu(B t^{-1}) - a,0 \big)
               \le \max \bigg( \sum_s g(st) \mu(s),0 \bigg)    \\
       \Rightarrow \quad g(t) \le \sum_s g(st) \mu(s).
\end{gathered}
\]
Thus
\[
     g(t) = \sum_s g(st) \mu(s), \quad t \in S(\nu).
\]
Now for $t \in A$, $g(t) = 0$
\[
\begin{aligned}
       \Rightarrow &\quad g(st) = 0 \hbox{ for } s \in S(\mu), \quad
                    t \in A \cap S(\nu)                        \\
       \Rightarrow &\quad \nu \big[ B(t^{-1}) s^{-1} \big]
               \le a, \quad t \in A \cap S(\nu), \ S \in S(\mu).
\end{aligned}
\]
Now $x \in A \cap S(\nu)$, $\nu(Bx^{-1}) = a$
\[
     \Rightarrow \quad \nu(Bx^{-1}s^{-1}) \le \nu(Bx^{-1})
          \hbox{ for all } s \in S(\mu).
\]
Now since $\nu = \nu * \mu$
\[
\begin{gathered}
       \nu(Bx^{-1}) = \sum_{s \in S(\mu)} \nu(Bx^{-1}y^{-1}) \mu(y) \\
       \sum_y \big[ \nu(Bx^{-1}) - \nu(Bx^{-1}y^{-1}) \big] u(y)
               = 0 \Rightarrow \nu(Bx^{-1}) = \nu(Bx^{-1}y^{-1}) \quad \qed
\end{gathered}
\]
\end{proof}

Now let $w_G$ be the uniform measure on a finite group $G$,
$|G| = n$.  Let $x \in G$ and consider the delta measure at $x,\dl_X$
and let $B \ub G$.  Then
\[
\begin{aligned}
       w_G * \dl_X(B) &= \sum w_G (Bh^{-1}) \dl_X(h)           \\
       &= w_G(Bx^{-1})                                         \\
       &= |B| \frac{1}{n}                                      \\
       &= w_G(B).
\end{aligned}
\]
We want to generalize this notion of translation invariance to
semigroups.

\begin{definition}   \label{defn:}
Let $\mu \in \Prob(S)$.  
Then $\mu$ is called $r^*$-invariant if for $B \ub S$, $\mu = \mu * \dl_X \big[ \mu(B) = \mu(Bx^{-1}) \big]$ if $\mu = \dl_X * \mu$.
\end{definition}


\runinhead{Homework.} Show that if $\mu$ is $r^*$-invariant or $1^*$-invariant, then $\mu^2 = \mu$.

Notice that if $\mu$ is $r^*$-invariant then
\[
     S(\mu) = S(\mu * \dl_X) = S(\mu) \cdot x.
\]
$1^*$-invariant
\[
     S(\mu) = S(\dl_X *\mu) = x \cdot S(\mu) .
\]

\begin{lemma} 		\label{lem:3.2} 
If $\mu$ is $r^*$-invariant, then for $x \in S(\mu)$, $x \cdot S(\mu)$ is left cancellative $(zx = zy \Rightarrow x = y)$.
\end{lemma}

\begin{proof} Let $\mu$ be a $r^*$-invariant measure on $S$.
Then for any $s \in S(\mu)$, $S(\mu) \cdot s = S(\mu)$.  
Take arbitrary elements in $x-S(\mu)$, $xy$, $xz$, $xw$.

Suppose $(xy)(xz) = (xy)(xw)$.  Now $S(\mu) \cdot (xyx) = S(\mu)$.  
Thus $t(xyx) = x$ for some $t \in S(\mu)$.  
Then,
\[
\begin{aligned}
       t(xyx)z &= t(xyx)w                                      \\
       \Rightarrow \quad xz &= xw
\end{aligned}
\]
Thus $x \cdot S(\mu)$ is left-cancellative.
\qed
\end{proof}

\runinhead{Homework.} Show that if $\mu$ is $1^*$-invariant $S(\mu) \cdot x$ is right-cancellative.

\begin{lemma}		\label{lem:3.3}  
Let $\mu \in \Prob(S)$ be both $r^*$ and $1^*$-invariant.  
Then $S(\mu)$ is a group and $\mu$ is the uniform measure on $S(\mu)$.
\end{lemma}

\begin{proof} Since $\mu = \mu^2$, then $S(\mu) = S(\mu)
     \cdot S(\mu)$ and $S(\mu)$ is a semigroup.  Now for
     $x \in S(\mu)$,
\[
\begin{aligned}
       S(\mu) \cdot x = S(\mu) &\hbox{ by } r^*\hbox{-invariance}  \\
       x \cdot S(\mu) = S(\mu) &\hbox{ by } 1^*\hbox{-invariance}.
\end{aligned}
\]
Thus $S(\mu)$ is both right and left simple and is thus a group.

Let $y,z \in S(\mu)$ and suppose $z \in yx^{-1}$ for $x \in S(\mu)$.  
Since $\mu = \mu * \dl_X$
\[
\begin{aligned}
       \mu(y) &= \mu * \dl_X(y) = \mu(yx^{-1})     \\
       &= \mu(yx^{-1} \cap S(\mu) ) = \mu(z)
\end{aligned}
\]
where $\{ z \} = yx^{-1} \cap S(\mu)$ is a single element since
$S(\mu)$ is a group.

Thus $\mu(y) = \mu(z)$ for any pair $y,z \in S(\mu)$ which show
$\mu$ is uniform on $S(\mu)$.
\qed
\end{proof}

\begin{theorem} 	\label{thm:3.1}
Suppose for $\mu \in \Prob(S)$, $\mu = \mu^2$.
Then $S(\mu)$ is a completely simple semigroup of $S$.
\end{theorem}

\begin{proof}  Now $S(\mu) \cdot S(\mu) = S(\mu)$ (since $\mu * \mu = \mu$). 
Thus
\[
     S(\mu) = \bigcup_{x \in S(\mu)} S(\mu) \cdot x.
\]
If we can show that each $S(\mu) \cdot x$ is a minimal left ideal,
then $S(\mu)$ will be completely simple.

Now, anytime $S \cdot S = S$, $S = \bigcup\limits_{x \in S}
S \cdot x$.  
But if $S$ is not completely simple, then for some $x \in S$, $S \cdot x$ is not minimal.

Now, to show that $S(\mu) \cdot x$ is left simple we need to show
that $(S(\mu) \cdot x) \cdot L = S(\mu) \cdot x$ for all subsets
$L \ub S(\mu) \cdot x$.  
Clearly, it is sufficient to show this for each element $z \in S(\mu) \cdot x$.

Call $\mu * \dl_x = \mu_x$ and let $z \in S(\mu) \cdot x$.  
Then $z = y \cdot x$ for some $y \in S(\mu)$.
\[
\begin{aligned}
       \mu * \dl_x * \dl_2(B) = \mu * \dl_x(Bz^{-1})
               = \mu_x(Bz^{-1}) &= \mu_x (B(y \cdot x)^{-1})   \\
       &= \mu_x B(x^{-1}y^{-1})                                \\
       &= \mu\, B(x^{-1}y^{-1}x^{-1})                           \\
       &= \mu \big[ Bx^{-1} (y)^{-1} \big]                     \\
\noalign{by $xy \in S(\mu)$}
       &= \mu(Bx^{-1}) = \mu_x(B) = \mu * \dl_x(B).
\end{aligned}
\]
Thus
\[
\begin{alignedat}{2}
       & &\quad
       \mu * \dl_x * \dl_z &= \mu * \dl_x                      \\
       &\Rightarrow &\quad S(\mu) * \dl_x *\dl_z &= S(\mu * \dl_x)       \\
       &\Rightarrow &\quad [S(\mu) \cdot x] \cdot z &= S(\mu) \cdot x    \\
\end{alignedat}
\]
for all $z \in S(\mu) \cdot x$
\[
  \Rightarrow \quad S(\mu) \cdot x
	\hbox{ is left simple.} 
\]
Thus each $S(\mu) \cdot x$ is a minimal left ideal and
\[
     S(\mu) = \bigcup_{x \in S(\mu)} S(\mu) \cdot x
\]
is a union of minimal left ideals and is completely simple.
\qed
\end{proof}

\begin{theorem} 	\label{thm:3.2}
Let $\mu^2=\mu$ and let $S(\mu) \cong x \times G \times Y$.
Then there exists $\mu_1 \in \Prob(X)$, $\mu_2 \in \Prob(G)$, $\mu_3 \in \Prob(Y)$ such that $\mu = \mu_1 * \mu_2 * \mu_3$ where $\mu_2$ is the uniform measure on $G$.  
Conversely, if $\mu_1,\mu_3 \in \Prob(S)$ and $\mu_2$ is the uniform measure on a subgroup $G$ of $S$, with $S(\mu_3) \cdot S(\mu_1) \ub G$, then $\mu_1 * \mu_2 * \mu_3$ is idempotent.
\end{theorem}

\begin{proof} Now, by the previous theorem, $S(\mu)$ is completely simple.  Let $e = e^2 \in S(\mu)$.  
Consider $\dl_e * \mu * \dl_e$.  
By the previous theorem,
\[
     (\dl_e * \mu * \dl_e) * \dl_{ese} = \dl_e * \mu * \dl_e
\]
and
\[
     \dl_{ese} * (\dl_e * \mu * \dl_e) = \dl_e * \mu * \dl_e.
\]
Thus $\dl_e * \mu * \dl_e$ is both $r^*$- and $\ell^*$-invariant, and
thus is the uniform measure on its support, $e \cdot S(\mu) \cdot e$.

Also, notice that for $B \ub S$,
\[
\begin{aligned}
       \mu * \dl_e * \mu(B) &= \sum \mu * \dl_e (Bx^{-1}) \mu(x) \\
       &= \sum \mu(Bx^{-1}e^{-1}) \mu(x)                       \\
       &= \sum \mu(Bx^{-1}) \mu(x)                             \\
       &= \mu * \mu(B) = \mu(B).
\end{aligned}
\]
Thus $\mu * \dl_e * \mu = \mu$ and $(\mu * \dl_e) * (\dl_e * \mu) =
\mu$.  
Now, let $X = E[S(\mu) \cdot e]$, $G = e \cdot S(\mu) \cdot e$,
$Y = E[e \cdot S(\mu)]$.

Define $\mu_1(A) = \mu * \dl_e \left( \Phi \big[ A \times G \times \{ e \} \big] \right)$.  
Notice $\mu_1 \in \Prob(X)$ and $\mu_2(B) = \mu * \dl_e \left( \Phi \big[ X \times B \times \{ e \} \big] \right)$.  
Now $A \ub X = E(S(\mu) \cdot e)$ and $S(\mu) = S(\mu * \dl_e) = S(\mu) \cdot e$.
Also $B \ub G = eS(\mu)e \ub E [S(\mu) \cdot e] \cdot eS(\mu) e \ub
S(\mu) \cdot e$.

Thus both $\mu_1,\mu_2$ are $r^*$-invariant, but since $S(\mu_2) = G$, then $\mu_2$ must be the Haar measure on $G$.  
If $A \ub X$, $B \ub G$,
\[
\begin{aligned}
       \mu_1 * \mu_2(AB) = \mu_1 * \mu_2(A \cdot B \cdot e)
               &= \sum_{y\in G} \mu_1
                    \big( \{ ABe \} y^{-1} \big) \mu_2(y)      \\
       &= \sum_{y\in G} \mu_1 \big[ \{ ABe \} y^{-1} \cap X
                    \big] \mu_2(y),
\end{aligned}
\]
but $\{ ABe \} y^{-1} = \{ z \mid zy \in ABe, \ y \in G \}$.  
If $zy \in ABe$, $z \in X$, $y \in G$, then $zye \in ABe$.  However, recall that $\Phi: X \times G \times Y \to X \cdot G \cdot Y$ is an isomorphism.
Since $zye \in X \cdot G \cdot Y$
\[
     \Phi^{-1} (zye) = (z,y,e)
\]
and for each $y \in B$, $\big[ \{ ABe \} y^{-1} \big] = A$ and
otherwise $[ABe] y^{-1} = \emptyset$.  
Thus
\[
\begin{aligned}
       \mu_1 * \mu_2(AB) &= \sum_{y \in G} \mu_1
               \big[ \{ ABe \} y^{-1} \big] \mu_2(y)           \\
       &= \sum_{y \in B} \big[ \mu_1\{ ABe \} y^{-1} \big] \mu_2(y) \\
       &\qquad + \sum_{y \in G-B}
               \big( \mu_1\{ ABe \} y^{-1} \big) \mu_2(y)           \\
       &= \mu_1(A) \sum_{y \in B} \mu_2(y)                     \\
       &= \mu_1(A) \mu_2(B).
\end{aligned}
\]

Similarly, if $\mu_3(C) = \dl_e * \mu \big( \Phi[e \times G \times C]
\big)$ for $C \ub E[e \cdot S(\mu)] = Y$, $\mu_2 * \mu_3(BC) = \mu_2(B) \mu_3(C)$.

Since $\mu * \dl_e$ is $r^*$-invariant, it is idempotent and $(\mu * \dl_e) * (\mu * \dl_e) = \mu * \dl_e$.  
Thus if $A,B$ are as above,
\[
\begin{aligned}
       \mu * \dl_e \big( \Phi[A \times B] \big) &= (\mu * \dl_e)
               * (\mu * \dl_e) \big[ \Phi(A \times B) \big]        \\
       &= \sum_{x \in S(\mu) \cdot e}
               \mu * \dl_e \big( x^{-1}
               \Phi [A \times B \times \{ e \} ] \big) \mu * \dl_e(x).
\end{aligned}
\]
Now
\[
     \Phi^{-1} [S(\mu) \cdot e] = \Phi^{-1} [XG(Y \cdot e) \cdot e]
          = X \times G \times \{ e \}.
\]
Thus
\[
     x \in S(\mu) \cdot e \Rightarrow \Phi^{-1}(x) \in X \times G \times \{ e \}.
\]
Also
\[
     X^{-1} \Phi[A \times B \times \{ e \}]
          = \{ z \mid xz \in \Phi[A \times B \times e] \}
\]
\[
\begin{aligned}
       &\Rightarrow \quad \Phi^{-1}(xz) \in A \times B \times \{ e \}          \\
       &\Rightarrow \quad \Phi^{-1}(x) \Phi^{-1}(z) \in A \times B \times \{ e \}  \\
       &\Rightarrow \quad (x_1,g_1,e)(z_1,z_2,z_3)
                    \in A \times B \times \{ e \}                      \\
       &\Rightarrow \quad (x_1,g_1,(ez_1) z_2,z_3)
                    \in A \times B \times \{ e \}                      \\
       &\Rightarrow \quad x_1 \in A, \ \ z_2 \in g^{-1}_1 B.
\end{aligned}
\]

     Finally,
\[
\begin{aligned}
       \mu * \dl_e \big( \Phi[A \times B] \big)
          &= \sum_{x \in S(\mu) \cdot e} \mu * \dl_e
               \big( x^{-1} \Phi[A \times B \times \{ e \}] \big)
               \mu * \dl_e(x)                                  \\
          &= \sum_{(x_1,g_1,e) \in A \times G \times \{ e \}}
               \mu * \dl_e
               \big( \Phi[ X \times g^{-1}_1B \times \{ e \}] \big)
               \mu * \dl_e(x) \big[ \Phi(x_1,g_1,e) \big]      \\
          &= \sum_{(x_1,g_1,e) \in A \times G \times \{ e \}}
               \mu_2 (g^{-1}_1B) \mu * \dl_e \big[ \Phi
               (x_1,g_1,e) \big]                               \\
          &= \mu_2(B) \sum_{(x_1,g_1,e) \in A \times G \times \{ e \}}
               \mu * \dl_e \big[ \Phi(x_1,g_1,e) \big]         \\
          &= \mu_1(A) \mu_2(B).
\end{aligned}
\]

Now since $\mu_1 * \mu_2(AB) = \mu_1(A) \mu_2(B)$ and $\mu * \dl_e(AB) = \mu_1(A) \mu_2(B)$.  Then $\mu_1 * \mu_2 = \mu * \dl_e$.

Similarly, $\mu_2 * \mu_3 = \dl_e * \mu$.
\[
  (\mu_1 * \mu_2 * \mu_3) = (\mu_1 * \mu_2) * (\mu_2 * \mu_3)
       = (\mu * \dl_3) * (\dl_e * \mu) = \mu.
\quad \qed
\]
\end{proof}

\runinhead{Midterm.}
\[
     (\mu_1 * \mu_2 * \mu_3(ABC)) = \mu_1(A) * \mu_2(B) * \mu_3(C)
\]
